%%%%%%%%%%  scan idx 195  |  folio 182  %%%%%%%%%%
dont l'image essentielle est $\mathrm D^b(A)_{\coh}$ (sous-catégorie triangulée de $\mathrm D(A)$ formée des complexes à cohomologie bornée et pseudo-cohérents relativement à la catégorie des modules libres de type fini).
\item Soit $E\in\Ob\mathrm D(S)$, et soit $[m,n]$ un intervalle de $\mathbb Z$. Si $E$ est d'amplitude parfaite contenue dans $[m,n]$ (I 4.8), $E$ est isomorphe à un complexe $F$ tel que $F^i=0$ pour $i\in[m,n]$, $F^i\in\Ob\Lib(S)$ pour $i>m$ et $F^m\in\Ob\Loclib(S)$.
\item Le foncteur canonique
\[
\mathrm K^b(\Loclib(S))\longrightarrow\mathrm D(S)_{\parf}\qquad(\text{I }4.9)
\]
est une équivalence de catégories. En outre, le foncteur $\mathbb R\Gamma_S:\mathrm D^+(S)\to\mathrm D^+(A)$ induit un foncteur pleinement fidèle
\[
\mathbb R\Gamma_S:\mathrm D(S)_{\parf}\longrightarrow\mathrm D^+(A),
\]
dont l'image essentielle est $\mathrm D(A)_{\parf}$ (sous-catégorie triangulée de $\mathrm D(A)$ formée des complexes parfaits relativement à la catégorie des modules projectifs de type fini).
\end{enumerate}

\emph{Preuve.} a) Les foncteurs
\[
\mathrm K^{-,b}(\Lib(S))\longrightarrow\mathrm K^{-,b}(\Loclib(S))\longrightarrow\mathrm D^b(\Coh(S))\xrightarrow{\Gamma_S}\mathrm D^b(A)_{\coh}
\]
sont des équivalences de catégories en vertu de (2.4.3 b) et c)). Comme le foncteur (pleinement fidèle) $\Gamma_S:\mathrm D^b(\Coh(S))\to\mathrm D(A)$ se factorise en
\[
\mathrm D^b(\Coh(S))\longrightarrow\mathrm D^b(S)_{\coh}\xrightarrow{\mathbb R\Gamma_S}\mathrm D(A),
\]
on conclut grâce à (2.4.4).

b) Supposons $\operatorname{parf.amp}(E)\subset[m,n]$. En particulier, $E$ est pseudo-cohérent ; en vertu de a), on peut donc supposer, quitte à remplacer $E$ par un objet isomorphe, que $E^i\in\Ob\Lib(S)$ pour $i>m$ et $E^i=0$ pour $i>n$.
